\subsection{由上闭链定义的主纤维丛}

定义 6. —— 设 B 为拓扑空间，G 为拓扑群，\(\mathcal{U} = (\mathrm{U}_{i})_{i\in \mathrm{I}}\) 为 B 的开覆盖。B 上取值于 G、从属于 \(\mathcal{U}\) 的连续 1-上闭链，是指对 I 的元素的每个偶 \((i,j)\) ，给定一个连续映射 \(g_{i,j}\) （从 \(\mathrm{U}_i\cap \mathrm{U}_j\) 到 G 中），使得对每个三元组 \((i,j,k)\in \mathrm{I}\times \mathrm{I}\times \mathrm{I}\) 以及 \(\mathrm{U}_i\cap \mathrm{U}_j\cap \mathrm{U}_k\) 的每个点 b，有

\[
g _ {i, k} (b) = g _ {i, j} (b) g _ {j, k} (b).
\]

我们以 \(Z_{\mathrm{cont}}^{1}(\mathrm{B},\mathcal{U},\mathrm{G})\) 表示 B 上取值于 G、从属于覆盖 U 的连续 1-上闭链组成的集合。

在本分册中，我们把连续 1-上闭链简称为上闭链。

设 \((\mathrm{E},p)\) 为 B 上的主 G-丛，\(\mathcal{U} = (\mathrm{U}_{i})_{i\in\mathrm{I}}\) 为由开集 \(U_{i}\) 组成的 B 的覆盖。若对每个 \(i\in I\)，E 在 \(U_{i}\) 上都可平凡化，则称 E 在 U 上可平凡化。此时我们把 E 的 U-平凡化称为族 \((f_{i})_{i\in\mathrm{I}}\) ，其中 \(f_{i}\colon p^{-1}(\mathrm{U}_{i})\to\mathrm{U}_{i}\times\mathrm{G}\) 是 E 在 \(U_{i}\) 上的一个平凡化。

设 \(\mathcal{U} = (\mathrm{U}_i)_{i\in \mathrm{I}}\) 为由开集组成的 B 的覆盖，\((\mathrm{E},p)\) 为 B 上赋予了 \(\mathcal{U}\)-平凡化 \((f_i)_{i\in \mathrm{I}}\) 的主 G-丛。对每个偶 \((i,j)\in \mathrm{I}\times \mathrm{I}\) ，以 \(g_{ij}\) 表示从 \(\mathrm{U}_i\cap \mathrm{U}_j\) 到 G 中由

\[
(b, g _ {i j} (b)) = f _ {i} \circ f _ {j} ^ {- 1} (b, e),
\]

所定义的映射，其中 e 表示 G 的单位元。它是连续的。

由于 \(f_{i}\) 与 \(f_{j}\) 与 G 的作用相容，我们有

\[
(f _ {i} \circ f _ {j} ^ {- 1}) (b, g) = (b, g _ {i j} (b) g)
\]

其中 \(b \in \mathrm{U}_i \cap \mathrm{U}_j\) ，\(g \in \mathrm{G}\) 。若 \(b\) 是 \(\mathrm{U}_i \cap \mathrm{U}_j \cap \mathrm{U}_k\) 的点（其中 \(i, j, k \in \mathrm{I}\)），则有

\[
(f _ {i} \circ f _ {k} ^ {- 1}) (b, e) = (b, g _ {i k} (b))
\]

以及

\[
\begin{array}{c} (f _ {i} \circ f _ {k} ^ {- 1}) (b, e) = (f _ {i} \circ f _ {j} ^ {- 1}) \circ (f _ {j} \circ f _ {k} ^ {- 1}) (b, e) \\ = (f _ {i} \circ f _ {j} ^ {- 1}) (b, g _ {j k} (b)) = (b, g _ {i j} (b) g _ {j k} (b)). \end{array}
\]

由此得出

\[
g _ {i k} (b) = g _ {i j} (b) g _ {j k} (b),
\]

故族 \((g_{ij})\) 是 B 上取值于 G、从属于覆盖 U 的上闭链。它称为由平凡化族 \((f_{i})_{i\in\mathrm{I}}\) 定义的上闭链。

设 B、G 与 U 如定义 6 所述，\((g_{ij}) \in Z_{\mathrm{cont}}^{1}(\mathrm{B}, \mathcal{U}, \mathrm{G})\) 为一个上闭链。对每个偶 \((i, j) \in \mathrm{I} \times \mathrm{I}\) ，映射 \(\gamma_{ij} \colon (\mathrm{U}_{i} \cap \mathrm{U}_{j}) \times \mathrm{G} \to (\mathrm{U}_{i} \cap \mathrm{U}_{j}) \times \mathrm{G}\) 由

\[
\gamma_ {i j} (b, g) = \left(b, g _ {i j} (b) g\right) \quad \text {   for   } b \in \mathrm{U} _ {i} \cap \mathrm{U} _ {j} \text {   and   } g \in \mathrm{G}\tag{5}
\]

所定义，它是基 \(U_{i} \cap U_{j}\) 上主丛的同构。对每个三元组 \((i,j,k) \in \mathrm{I} \times \mathrm{I} \times \mathrm{I}\) 以及每个偶 \((b,g) \in (\mathrm{U}_{i} \cap \mathrm{U}_{j} \cap \mathrm{U}_{k}) \times \mathrm{G}\) ，我们有：

\[
\gamma_ {i k} (b, g) = \gamma_ {i j} \circ \gamma_ {j k} (b, g).
\]

设 F 为把空间 \(U_{i} \times G\) 沿 \((\mathrm{U}_{i} \cap \mathrm{U}_{j}) \times \mathrm{G}\) 借助双射 \(\gamma_{ij}\) 粘合起来所得的拓扑空间（TG, I, 第 16 页）。对每个 \(i \in I\)，\(U_{i} \times G\) 在 F 中的像是 F 的开集（TG, I, 第 17 页, 命题 9）。典范投影 \(p_{i}: U_{i} \times G \to U_{i}\) 粘合成一个从 F 到 B 的连续映射 p。由于诸映射 \(\gamma_{ij}\) 与空间 \(U_{i} \times G\) 中的 G 右作用相容，由此导出 G 在 F 上的连续右作用，它使 F 成为以 B 为基、以 G 为群、且在每个 \(U_{i}\) 上带有平凡化的主纤维丛。我们称 F 为由上闭链 \((g_{ij})\) 定义的主纤维丛。

定义 7. —— 设 B 为拓扑空间，G 为拓扑群，\(\mathcal{U} = (\mathrm{U}_{i})_{i\in \mathrm{I}}\) 为 B 的开覆盖。称两个上闭链 \((g_{ij})\) 与 \((g_{ij}^{\prime})\) （均属于 \(\mathrm{Z}_{\mathrm{cont}}^{1}(\mathrm{B},\mathcal{U},\mathrm{G})\) ）是上同调的，如果存在族 \((h_{i})_{i\in\mathrm{I}}\) （由连续映射 \(h_{i}\colon U_{i}\to G\) 组成），使得

\[
g _ {i j} ^ {\prime} (b) = h _ {i} (b) g _ {i j} (b) h _ {j} (b) ^ {- 1}\tag{6}
\]

对每个偶 \((i,j)\in \mathrm{I}\times \mathrm{I}\) 以及每个 \(b\in \mathrm{U}_i\cap \mathrm{U}_j\) 成立。

关系"\((g_{ij})\) 与 \((g'_{ij})\) 上同调"是集合 \(Z_{\mathrm{cont}}^{1}(\mathrm{B},\mathcal{U},\mathrm{G})\) 上的一个等价关系。我们以 \(H_{\mathrm{cont}}^{1}(\mathrm{B},\mathcal{U},\mathrm{G})\) 表示 \(Z_{\mathrm{cont}}^{1}(\mathrm{B},\mathcal{U},\mathrm{G})\) 关于此等价关系的商集。

命题 11. —— 设 B 为拓扑空间，G 为拓扑群，\(\mathcal{U} = (\mathrm{U}_{i})_{i\in \mathrm{I}}\) 为 B 的开覆盖。

\begin{enumerate}
\def\labelenumi{\alph{enumi})}
\item
  B 上每个在 U 上可平凡化的主 G-丛都同构于由 \(Z_{\mathrm{cont}}^{1}(\mathrm{B}, \mathcal{U}, \mathrm{G})\) 中某个上闭链定义的主丛。
\item
  设 (E, p) 与 (E', p') 为 B 上在 U 上可平凡化的主 G-丛。设 \((f_i)_{i \in I}\) （相应地 \((f_i')_{i \in I}\) ）为 (E, p) （相应地 \((E', p')\) ）的适合 U 的平凡化，并以 \((g_{ij})_{(i,j) \in I \times I}\) （相应地 \((g_{i,j}')_{(i,j) \in I \times I}\) ）表示由此平凡化定义的上闭链。此时，主纤维丛 (E, p) 与 (E', p') 同构，当且仅当这些上闭链上同调。
\end{enumerate}

我们来证明 \(b\) 。设 \(\varphi: \mathrm{E} \to \mathrm{E}'\) 为以 B 为基、以 G 为群的主丛的同构。对每个 \(i \in \mathrm{I}\) ，设 \(h_i\) 为从 \(\mathrm{U}_i\) 到 G 中的连续映射，由

\[
(b, h _ {i} (b)) = f _ {i} ^ {\prime} \circ \varphi \circ f _ {i} ^ {- 1} (b, e).
\]

所定义。由于对每个 \(i \in I\) ，\(f_{i}\) 与 \(f_{i}^{\prime}\) 都与 G 的作用相容，故对每个 \(b \in U_{i}\) 与每个 \(g \in G\) ，有

\[
(f _ {i} ^ {\prime} \circ \varphi \circ f _ {i} ^ {- 1}) (b, g) = (b, h _ {i} (b) g)
\]

以及

\[
(f _ {i} \circ \varphi^ {- 1} \circ (f _ {i} ^ {\prime}) ^ {- 1}) (b, g) = (b, h _ {i} (b) ^ {- 1} g).
\]

因此，对每个偶 \((i,j)\in \mathrm{I}\times \mathrm{I}\) 以及每个点 \(b\) （属于 \(\mathrm{U}_i\cap \mathrm{U}_j\) ），

\[
\begin{array}{c} f _ {i} ^ {\prime} \circ (f _ {j} ^ {\prime}) ^ {- 1} (b, e) = (f _ {i} ^ {\prime} \circ \varphi \circ f _ {i} ^ {- 1}) \circ (f _ {i} \circ f _ {j} ^ {- 1}) \circ (f _ {j} \circ \varphi^ {- 1} \circ (f _ {j} ^ {\prime}) ^ {- 1}) (b, e) \\ = (b, h _ {i} (b) g _ {i j} (b) h _ {j} (b) ^ {- 1}) \end{array}
\]

从而有

\[
g _ {i j} ^ {\prime} (b) = h _ {i} (b) g _ {i j} (b) h _ {j} (b) ^ {- 1}.
\]

这就证明上闭链 \((g_{ij})\) 与 \((g'_{ij})\) 上同调。

反之，设这些上闭链上同调，并设 \((h_{i})_{i\in\mathrm{I}}\) 为由连续映射 \(h_{i}\colon U_{i}\to G\) 组成的族，它满足 \(g_{ij}^{\prime}(b)=h_{i}(b)g_{ij}(b)h_{j}(b)^{-1}\) （其中 \(i,j\in I\)，\(b\in U_{i}\cap U_{j}\)）。对 \(i\in I\) ，设 \(\varphi_{i}\colon p^{-1}(U_{i})\to(p')^{-1}(U_{i})\) 为由 \(f_{i}^{\prime}\circ\varphi_{i}\circ f_{i}^{-1}(b,g)=(b,h_{i}(b)g)\) （其中 \(b\in U_{i}\)，\(g\in G\)）定义的映射。这是以 \(U_{i}\) 为基、以 G 为群的主纤维丛的同构。对 \((i,j)\in I\times I\) ，分别以 \(\gamma_{ij}\) 与 \(\gamma_{ij}^{\prime}\) 表示如上依附于上闭链 \((g_{ij})\) 与 \((g_{ij}^{\prime})\) 的粘合同胚（I, 第 115 页, 公式 (5)），使得 \(f_{i}(b,g)=\gamma_{ij}\circ f_{j}(b,g)\) 与 \(f_{i}^{\prime}(b,g)=\gamma_{ij}^{\prime}\circ f_{j}^{\prime}(b,g)\) 对一切 \((b,g)\in(U_{i}\cap U_{j})\times G\) 成立。于是，对 \((i,j)\in I\times I\) 与 \((b,g)\in(U_{i}\cap U_{j})\times G\) ，有关系

\[
\begin{array}{r l} & f _ {i} ^ {\prime} \circ \varphi_ {j} \circ f _ {i} ^ {- 1} (b, g) = \gamma_ {i j} ^ {\prime} \circ (f _ {j} ^ {\prime} \circ \varphi_ {j} \circ f _ {j} ^ {- 1}) (b, g _ {i j} (b) ^ {- 1} g) \\ & \qquad = \gamma_ {i j} ^ {\prime} (b, h _ {j} (b) g _ {i j} (b) ^ {- 1} g) \\ & \qquad = (b, g _ {i j} ^ {\prime} (b) h _ {j} (b) g _ {i j} (b) ^ {- 1} g) \\ & \qquad = (b, h _ {i} (b) g) = f _ {i} ^ {\prime} \circ \varphi_ {i} \circ f _ {i} ^ {- 1} (b, g). \end{array}
\]

这表明 \(\varphi_{i}\) 与 \(\varphi_{j}\) 在 \(p^{-1}(\mathrm{U}_{i}\cap\mathrm{U}_{j})\) 上相合。于是诸态射 \(\varphi_{i}\) 粘合成一个从 E 到 E' 的主丛的 B-态射。断言 \(b\) 由"任何以 B 为基、以 G 为群的主丛态射都是同构"得出（I, 第 93 页, 命题 1）。

现在来证明 \(a\) 。设 \((\mathrm{E}, p)\) 为以 G 为结构群的主纤维丛，且对每个 \(i \in \mathrm{I}\) ，带有平凡化 \(f_i: p^{-1}(\mathrm{U}_i) \to \mathrm{U}_i \times \mathrm{G}\) （在 \(\mathrm{U}_i\) 上）。设 \((g_{ij})\) 为由此族定义的上闭链，F 为由上闭链 \((g_{ij})\) 定义的主纤维丛。按构造，主纤维丛 F 带有 \(\mathcal{U}\) 上的一个平凡化；该平凡化定义出上闭链 \((g_{ij})\)。由断言 \(b\) ，主纤维丛 \((\mathrm{E}, p)\) 同构于 F。

设 B 为拓扑空间，G 为拓扑群。设 \((E,p)\) 为以 B 为基、以 G 为群的主纤维丛。设 s 为满射 p 的一个截面（E, II, 第 19 页, 命题 8）。映射 \(f\colon B\times G\to E\) 由 \(f(b,g)=s(b)\cdot g\) 所定义，它是与 G 的作用相容的双射。给 \(B\times G\) 赋予经结构转移得到的拓扑；此时 \((B\times G,\mathrm{pr}_{1})\) 是以 B 为基、以 G 为群且同构于 E 的主纤维丛。因此，存在一个以 B 为基、以 G 为群的主丛的集合 T，使得每个以 B 为基、以 G 为群的主纤维丛都同构于 T 的某个元素。以 P(B, G) 表示以 B 为基、以 G 为群的主丛的同构类组成的集合（E, II, 第 47 页）。

设 \(\mathcal{U} = (\mathrm{U}_i)_{i\in \mathrm{I}}\) 为 B 的开覆盖。以 \(\mathrm{P(B,\mathcal{U},G)}\) 表示 \(\mathrm{P(B,G)}\) 中由在 \(\mathcal{U}\) 上可平凡化的主丛的同构类组成的子集。由命题 11，存在映射 \(r_{\mathcal{U}}\) （从 \(\mathrm{H}_{\mathrm{cont}}^{1}(\mathrm{B},\mathcal{U},\mathrm{G})\) 到 \(\mathrm{P(B,\mathcal{U},G)}\) 中），它把上闭链 \((g_{ij})\in Z_{\mathrm{cont}}^{1}(\mathrm{B},\mathcal{U},\mathrm{G})\) 的类对应到由此上闭链定义的主丛的同构类；它是双射（同前引）。逆双射 \(s_{\mathcal{U}}\) 把 \(\mathrm{P(B,\mathcal{U},G)}\) 中某个在 \(\mathcal{U}\) 上可平凡化的主纤维丛 E 的同构类，对应到由 E 在 \(\mathcal{U}\) 上任一平凡化族所定义的上闭链的类。在此对应下，平凡主纤维丛 \(\mathrm{B}\times \mathrm{G}\) 的同构类对应于常值上闭链 \(g_{ij} = e\) 的上同调类，后者也称为平凡上闭链。

依据主纤维丛的定义，集合 P(B, G) 是形如 P(B, U, G) 的集合之并，其中 U 取遍 B 的开覆盖。

设 \(\mathcal{U} = (\mathrm{U}_i)_{i\in \mathrm{I}}\) 为 B 的开覆盖，\(\mathcal{V} = (\mathrm{V}_k)_{k\in \mathrm{K}}\) 为比 \(\mathcal{U}\) 更细的 B 的开覆盖。我们有 \(\mathrm{P(B,\mathcal{U},G)}\subset\) \(\mathrm{P(B,\mathcal{V},G)}\) ；以 \(i_{\mathcal{V}\mathcal{U}}\) 表示由此包含关系定义的典范单射。选取一个映射 \(\varphi \colon \mathrm{K}\to \mathrm{I}\) ，使得 \(\mathrm{V}_k\subset \mathrm{U}_{\varphi (k)}\) 对一切 \(k\in \mathrm{K}\) 成立。给定上闭链 \((g_{ij})\in \mathrm{Z}_{\mathrm{cont}}^1 (\mathrm{B},\mathcal{U},\mathrm{G})\) ，对每个偶 \((k,\ell)\in \mathrm{K}\times \mathrm{K}\) ，置 \(\overline{g}_{k\ell} = g_{\varphi (k)\varphi (\ell)}\mid \mathrm{V}_k\cap \mathrm{V}_\ell\) 。

族 \((\overline{g}_{k\ell})\) 是 B 上取值于 G、从属于 V 的上闭链。若 \((g'_{ij}) \in \mathrm{Z}_{\mathrm{cont}}^{1}(\mathrm{B}, \mathcal{U}, \mathrm{G})\) 是与上闭链 \((g_{ij})\) 上同调的上闭链，则上闭链 \((\overline{g}'_{k\ell})\) （由 \((g'_{ij})\) 得出）与上闭链 \((\overline{g}_{k\ell})\) 上同调。这给出映射

\[
c (\varphi) \colon \mathrm{H} _ {\mathrm{cont}} ^ {1} (\mathrm{B}, \mathcal {U}, \mathrm{G}) \to \mathrm{H} _ {\mathrm{cont}} ^ {1} (\mathrm{B}, \mathcal {V}, \mathrm{G})
\]

它把 \((g_{ij})\) 的类对应到 \((\overline{g}_{k\ell})\) 的类。

设 E 为以 B 为基、以 G 为群的主纤维丛，且对每个 \(i \in I\) ，设 \(f_i: E_{U_i} \to U_i \times G\) 为 \(U_i\)-主纤维丛 \(E_{U_i}\) 的一个平凡化。对每个 \(k \in K\) ，设 \(f_k' : E_{V_k} \to V_k \times G\) 为由 \(f_{\varphi(k)}\) 限制到子集而得的平凡化。设 \((g_{ij}) \in Z_{\mathrm{cont}}^1(\mathrm{B}, \mathcal{U}, \mathrm{G})\) 为由族 \((f_i)\) 定义的上闭链；由族 \((f_{k}^{\prime})\) 定义的上闭链正是上面定义的上闭链 \((\overline{g}_{k\ell})\) 。于是，下面的图式是交换的：

\[
\begin{array}{c} \mathrm{H} _ {\mathrm{cont}} ^ {1} (\mathrm{B}, \mathcal {U}, \mathrm{G}) ^ {c (\varphi)} \longrightarrow \mathrm{H} _ {\mathrm{cont}} ^ {1} (\mathrm{B}, \mathcal {V}, \mathrm{G}) \\ \Big \downarrow^ {r _ {\mathcal {U}}} \qquad \qquad \qquad \qquad \Big \downarrow^ {r _ {\mathcal {V}}} \\ \mathrm{P} (\mathrm{B}, \mathcal {U}, \mathrm{G}) ^ {i _ {\mathcal {V} \mathcal {U}}} \longrightarrow \mathrm{P} (\mathrm{B}, \mathcal {V}, \mathrm{G}). \end{array}
\]

由于映射 \(r_{\mathcal{U}}\) 与 \(r_{\mathcal{V}}\) 都是双射，映射 \(c(\varphi)\) 与 \(i_{\mathcal{V}\mathcal{U}}\) 一样是单射，且不依赖于 \(\varphi\) 的选取。此后我们将以 \(c_{\mathcal{V}\mathcal{U}}\) 代替 \(c(\varphi)\) 来记写。

设 \(\mathcal{R}\) 为 \(\mathfrak{P}(\mathfrak{P}(B))\) 中构成 B 的开覆盖的元素组成的集合。对 B 的每个开覆盖 \(\mathcal{U}\) ，存在开覆盖 \(\mathcal{V}\) （属于 \(\mathcal{R}\) ），使得 \(\mathcal{V}\) 比 \(\mathcal{U}\) 既更细又更粗。集合 \(\mathcal{R}\) 对于关系 \(\leqslant\) （由"\(\mathcal{U} \leqslant \mathcal{V}\) 当且仅当 \(\mathcal{V}\) 是比 \(\mathcal{U}\) 更细的覆盖"定义）是有序的且是滤子的。由前所述，我们已定义了归纳系 \((\mathrm{H}_{\mathrm{cont}}^{1}(\mathrm{B},\mathcal{U},\mathrm{G}),c_{\mathcal{V}\mathcal{U}})\) （关于滤子化有序集 \(\mathcal{R}\) ），且族 \((r_{\mathcal{U}})\) 是从 \((\mathrm{H}_{\mathrm{cont}}^{1}(\mathrm{B},\mathcal{U},\mathrm{G}),c_{\mathcal{V}\mathcal{U}})\) 到 \((\mathrm{P}(\mathrm{B},\mathcal{U},\mathrm{G}),i_{\mathcal{V}\mathcal{U}})\) 的双射映射的归纳系。若以 \(\mathrm{H}_{\mathrm{cont}}^{1}(\mathrm{B},\mathrm{G})\) 表示系 \((\mathrm{H}_{\mathrm{cont}}^{1}(\mathrm{B},\mathcal{U},\mathrm{G}),c_{\mathcal{V}\mathcal{U}})\) 的归纳极限，以 \(r\colon \mathrm{H}_{\mathrm{cont}}^{1}(\mathrm{B},\mathrm{G})\to \mathrm{P}(\mathrm{B},\mathrm{G})\) 表示族 \((r_{\mathcal{U}})\) 的归纳极限，则有：

定理 4. —— 映射 \(r\colon \mathrm{H}_{\mathrm{cont}}^{1}(\mathrm{B},\mathrm{G})\to \mathrm{P}(\mathrm{B},\mathrm{G})\) 是双射。

设 \(\mathcal{U} = (\mathrm{U}_i)_{i\in \mathrm{I}}\) 为 B 的开覆盖；以 \(c_{\mathcal{U}}\) 表示典范映射 \(\mathrm{H}_{\mathrm{cont}}^{1}(\mathrm{B},\mathcal{U},\mathrm{G})\to \mathrm{H}_{\mathrm{cont}}^{1}(\mathrm{B},\mathrm{G})\) 。若 \((g_{ij})\) 是 B 上取值于 G、从属于开覆盖 \(\mathcal{U}\) 的上闭链，则元素 \(c_{\mathcal{U}}((g_{ij}))\) （属于 \(\mathrm{H}_{\mathrm{cont}}^{1}(\mathrm{B},\mathrm{G})\) ）称为上闭链 \((g_{ij})\) 的上同调类。

